% !Mode:: "TeX:UTF-8"
\chapter{总结与展望}

\section{研究总结}
本文围绕博物馆导览中“对话可持续、情境可对齐、引导可感知”这一核心问题，开展了面向博物馆导览的对话式数字人展示研究与设计实践。研究首先通过前期调研与文献梳理，明确了当前对话式数字人展示在连续讲解组织、现场情境响应与多模态引导表达方面的主要不足；在此基础上，提出并细化了 CCEG 框架，将展示设计能力归纳为对话组织、情境对齐与具身引导三个相互协同的维度。

基于框架要求，本文完成了系统级设计与实现。对话组织部分围绕展品知识结构、分层讲解策略与探索式提问引导展开，使讲解能够在导览过程中形成由概览到深入的连续推进；情境状态部分围绕空间位置、导览对象、用户参观状态与导览进程构建状态维护机制，为内容选择、策略调整与转场衔接提供依据；具身交互部分从人物形象与性格、动作状态与注意力引导、语音风格与语调节奏、空间及展品环境协同等方面进行设计，实现讲解语义与表达行为的一致映射。

结合实验与分析结果，本文所设计的对话式数字人导览原型系统在知识讲述支持、情境响应能力与注意力引导效果方面表现出相对于基线方式的改进趋势。研究表明，对话式数字人展示的有效性并不只取决于问答能力本身，更取决于对话结构、现场状态与具身表达三者的协同组织。上述结果为面向博物馆学习过程的对话式数字人展示设计提供了可落地的框架与实践依据。

\section{研究不足}
尽管本文完成了框架提出、系统实现与初步验证，仍存在若干局限。第一，实验样本规模和人群类型仍较有限，当前结果更适用于验证系统可行性和趋势性结论，后续仍需扩大样本并覆盖更多年龄、文化背景与参观经验层次的用户。第二，实验场景以有限展区与典型展品为主，尚未覆盖大规模、多主题、跨楼层的复杂导览任务，系统在长时段连续导览中的稳定性与迁移性仍需进一步验证。第三，部分评估指标主要依赖主观问卷与基础行为数据，尚缺少更细粒度的客观指标，如注意力轨迹、长期学习保持和多次复访行为等。第四，当前系统中的角色与表达策略仍以规则化配置为主，对个体差异的自适应能力有待增强。

\section{未来展望}
未来研究可从以下几个方向继续推进。第一，开展更大规模、跨场馆的长期实地评估，检验原型系统在真实运营条件下的稳定性、泛化能力与持续使用价值。第二，强化个性化导览机制，在保持导览主线稳定的前提下，根据用户兴趣、停留行为与交互历史实现更细粒度的动态适配。第三，进一步深化多模态协同设计，将语音、动作、视线、界面提示与空间感知统一到更完整的实时调度机制中，提升引导清晰度与交互自然性。第四，完善评估体系，在主观体验指标之外引入更系统的客观过程指标与学习结果指标，形成可参考的对话式数字人展示评估范式。第五，探索与博物馆教育活动、社群传播和无障碍服务的联动应用，拓展对话式数字人展示在公共文化服务中的实际价值。

总体而言，面向博物馆导览的对话式数字人展示研究仍处于快速发展阶段。本文工作提供了一个从问题识别、框架构建到系统落地与实验验证的完整路径。随着生成式人工智能、空间感知与多模态交互技术的进一步发展，对话式数字人展示有望从“可用”走向“好用”，并在真实参观场景中持续支持观众的理解建构与文化体验。
